\label{fig:ejercicio6-viga}
\end{minipage}
\end{center}
\end{solucion}

\begin{ejercicio}[7: Voladizo con carga triangular]
Un voladizo de longitud $L=\qty{4}{\metre}$ soporta una carga descendente:
\[
q(x)=9\left(1-\frac{x}{4}\right)\quad[\mathrm{kN/m}],
\]
máxima en el empotramiento y nula en el extremo libre.

Obtener $V(x)$ y $M(x)$.
\end{ejercicio}

\begin{solucion}
Carga total:
\[
Q=\frac12(9)(4)=\qty{18}{\kilo\newton}.
\]
Actúa a:
\[
\bar x=\frac43\ \mathrm{m}
\]
desde el empotramiento.

La reacción vertical es:
\[
R_A=\qty{18}{\kilo\newton}.
\]
El momento de reacción vale:
\[
M_A^{\text{reacción}}=18\frac43=\qty{24}{\kilo\newton\metre}.
\]
El valor \textbf{interno} del momento justo a la derecha del empotramiento, con el convenio de momento flector positivo del curso, es:
\[
M(0)=-24\ \mathrm{kN\,m}.
\]

Como
\[
\frac{\dd V}{\dd x}=-q(x),
\]
integramos:
\[
V(x)=18-\int_0^x9\left(1-\frac{s}{4}\right)\dd s,
\]
y queda
\[
\boxed{V(x)=18-9x+\frac98x^2}.
\]
Se comprueba $V(4)=0$.

Ahora:
\[
\frac{\dd M}{\dd x}=V.
\]
Integrando desde $x=0$:
\[
\boxed{M(x)=-24+18x-\frac92x^2+\frac38x^3}.
\]
Y $M(4)=0$.

\begin{tip}
En un extremo libre sin carga puntual ni momento aplicado:
\[
V(L)=0,\qquad M(L)=0.
\]
Estas condiciones son una comprobación excelente.
\end{tip}
\end{solucion}

\begin{ejercicio}[8: Reconstruir la carga a partir del momento]
En un tramo de viga:
\[
M(x)=12x-3x^2\quad[\mathrm{kN\,m}].
\]
Determinar el cortante y la carga distribuida.
\end{ejercicio}

\begin{solucion}
Sabemos:
\[
V(x)=\frac{\dd M}{\dd x}.
\]
Por tanto:
\[
\boxed{V(x)=12-6x\quad[\mathrm{kN}]}.
\]
A continuación:
\[
\frac{\dd V}{\dd x}=-q.
\]
Como $\dd V/\dd x=-6$, resulta:
\[
\boxed{q=\qty{6}{\kilo\newton\per\metre}}.
\]

\begin{idea}
Los diagramas contienen información física suficiente para reconstruir las cargas:
\flowbox{$M\xrightarrow{\text{derivar}}V\xrightarrow{\text{derivar y cambiar signo}}q$}
\end{idea}
\end{solucion}

\begin{ejercicio}[9: Viga con rótula interna]
Una viga está empotrada en $A$. Entre $A$ y una rótula interna $B$ hay $AB=\qty{3}{\metre}$. Entre $B$ y un apoyo móvil $C$ hay $BC=\qty{2}{\metre}$.

Sobre $AB$ actúa $q=\qty{4}{\kilo\newton\per\metre}$ descendente. En el centro de $BC$ actúa $P=\qty{6}{\kilo\newton}$ descendente.

Calcular las reacciones.
\end{ejercicio}

\begin{solucion}
La clave es comenzar por el tramo $BC$. La rótula $B$ no transmite momento.

El tramo $BC$ es una viga simplemente apoyada con carga centrada, de modo que:
\[
R_B^{(BC)}=R_C=\frac{P}{2}=\boxed{\qty{3}{\kilo\newton}}.
\]
Por acción y reacción, el tramo $AB$ recibe en $B$ una fuerza descendente de $\qty{3}{\kilo\newton}$.
